\documentclass[10pt,a4paper]{amsart}
\usepackage[margin=19mm]{geometry}
\usepackage{amsmath,amssymb,mathtools}
\usepackage[scr=rsfs,cal=euler]{mathalfa}
\usepackage{xfrac}
\usepackage[unicode,hidelinks,bookmarks=false]{hyperref}
\newcommand{\ie}{i.e.\ }
\newcommand{\cf}{cf.\ }
\newcommand{\resp}{respectively\ }
\newcommand{\Subsection}{\subsection}
\providecommand{\nobreakdash}{\nobreak\hbox{-}}
\setlength{\emergencystretch}{2em}
\multlinegap0pt
\usepackage{multirow}
\usepackage{mathrsfs}
\usepackage{appendix}
\usepackage{textcomp}
\usepackage{manyfoot}
\usepackage{booktabs}
\usepackage{algorithm}
\usepackage{algorithmicx}
\usepackage{algpseudocode}
\usepackage{listings}
\newtheorem{theorem}{Theorem}%  meant for continuous numbers
\newtheorem{proposition}[theorem]{Proposition}% 
\newtheorem{example}{Example}%
\newtheorem{remark}{Remark}%
\newtheorem{axiom}{Axiom}
\newtheorem{corollary}[theorem]{Corollary}
\newtheorem{lemma}[theorem]{Lemma}
\newtheorem{definition}{Definition}%
\begin{document}
\title[Gauge Geometry of Hodge Zero-Mode Transport in Parameter-Dep]{Gauge Geometry of Hodge Zero-Mode Transport in Parameter-Dependent Topological Data Analysis}
\author{Satoshi Kanno, Rei Nishimura, Hiroshi Yamauchi, Yoshi-aki Shimada}
\date{}
\maketitle
\noindent\emph{Source and licence.} Gauge Geometry of Hodge Zero-Mode Transport in Parameter-Dependent Topological Data Analysis. Version arXiv:2605.28326v1 (CC BY 4.0). This material is distributed under the Creative Commons Attribution 4.0 International licence, http://creativecommons.org/licenses/by/4.0/, as recorded in the arXiv OAI licence metadata for this exact version and shown on the abstract page https://arxiv.org/abs/2605.28326v1. Full rights notice: licenses/arxiv-2605.28326-CC-BY-4.0.txt.\par\smallskip
\noindent\textit{Editorial scope.} Counted unit: the original \texttt{theorem} environment \texttt{-} at source lines 1495--1526 together with its complete proof at lines 1528--1542; the delivered document keeps source lines 291--1543 so that every referenced label and every citation resolves inside it. Native editable source is the paper's own concatenated TeX; the AMS wrapper is editorial formatting only. No publisher class, upstream script or unrelated artwork is redistributed.
\section{Zero-Mode Bundles and Its Geometry}

In this section, we realize the family of zero-mode spaces of the ordinary combinatorial Hodge Laplacian as a family of subspaces inside a common Hilbert space, and show that it forms a vector bundle on a regular region. We then introduce a natural Berry-type connection on this zero-mode bundle and define the associated curvature and holonomy. Finally, we state that, on regular regions, the zero-mode projection and the curvature expressed by the projection formula are stable under perturbations of the operator family.

The Laplacian used in this section is not the persistent Laplacian, but the ordinary combinatorial Hodge Laplacian associated with the complex at each parameter point. Therefore, the construction in this section describes the transport geometry of instantaneous homology spaces at each scale and external parameter, rather than the transport geometry of fixed-birth persistent feature spaces.

% \subsection{Zero-mode spaces and the zero-mode bundle}

% We retain the notation of Section~2. Let \(A\subset\mathbb R\) be the filtration-parameter set, let \(\Lambda\) be the additional parameter space, and let
% \[
% \{K(d,\lambda)\}_{(d,\lambda)\in A\times\Lambda}
% \]
% be a parameter-dependent filtration. That is, for each fixed \(\lambda\in\Lambda\), we assume that the map \(d\mapsto K(d,\lambda)\) is a filtration. We also assume that each \(K(d,\lambda)\) is a finite simplicial complex on a fixed finite vertex set \(V\).

% Let \(K_{\max}\) denote the maximal simplicial complex on \(V\), and for each degree \(q\) set
% \[
% H_q:=C_q(K_{\max}).
% \]
% Here \(H_q\) is regarded as a finite-dimensional Hilbert space over \(\mathbb R\) or \(\mathbb C\), equipped with the standard inner product for which the oriented \(q\)-simplices of \(K_{\max}\) form an orthonormal basis. Then, for every \((d,\lambda)\in A\times\Lambda\),
% \[
% C_q(K(d,\lambda))\subset H_q
% \]
% is naturally embedded.

% For each \((d,\lambda)\in A\times\Lambda\), let
% \[
% \Delta_q^{\mathrm{Hodge},\mathrm{nat}}(d,\lambda)
% \]
% denote the ordinary \(q\)-th combinatorial Hodge Laplacian on \(K(d,\lambda)\). Using the boundary operator
% \[
% \partial_q(d,\lambda):C_q(K(d,\lambda))\to C_{q-1}(K(d,\lambda)),
% \]
% it is defined by
% \[
% \Delta_q^{\mathrm{Hodge},\mathrm{nat}}(d,\lambda)
% =
% \partial_q(d,\lambda)^\ast\partial_q(d,\lambda)
% +
% \partial_{q+1}(d,\lambda)\partial_{q+1}(d,\lambda)^\ast.
% \]
% Here \(\ast\) denotes the adjoint with respect to the standard inner product. This operator is self-adjoint on \(C_q(K(d,\lambda))\).

% To compare Laplacians at different parameter values, we extend this operator to a self-adjoint operator on the common Hilbert space \(H_q\). With respect to the orthogonal decomposition
% \[
% H_q=C_q(K(d,\lambda))\oplus C_q(K(d,\lambda))^\perp,
% \]
% we define
% \[
% \Delta_q^{\mathrm{Hodge}}(d,\lambda)
% :=
% \Delta_q^{\mathrm{Hodge},\mathrm{nat}}(d,\lambda)\oplus I,
% \]
% where \(I\) is the identity on the orthogonal complement \(C_q(K(d,\lambda))^\perp\). This extension is self-adjoint and satisfies
% \[
% \ker \Delta_q^{\mathrm{Hodge}}(d,\lambda)
% =
% \ker \Delta_q^{\mathrm{Hodge},\mathrm{nat}}(d,\lambda).
% \]

% \begin{definition}
% For each \((d,\lambda)\in A\times\Lambda\), define
% \[
% E_0(d,\lambda):=\ker \Delta_q^{\mathrm{Hodge}}(d,\lambda)\subset H_q.
% \]
% We call this subspace the zero-mode space.
% \end{definition}

% By the ordinary discrete Hodge theorem, this zero-mode space realizes the ordinary homology of the complex \(K(d,\lambda)\).

% \begin{proposition}
% For each \((d,\lambda)\in A\times\Lambda\), there is a natural isomorphism
% \[
% E_0(d,\lambda)\cong H_q(K(d,\lambda)).
% \]
% In particular,
% \[
% \dim E_0(d,\lambda)=\beta_q(d,\lambda).
% \]
% \end{proposition}

% \begin{proof}
% The extended operator \(\Delta_q^{\mathrm{Hodge}}(d,\lambda)\) acts as the identity on the orthogonal complement. Hence its zero modes agree with the zero modes of the original Hodge Laplacian on \(C_q(K(d,\lambda))\). That is,
% \[
% \ker \Delta_q^{\mathrm{Hodge}}(d,\lambda)
% =
% \ker \Delta_q^{\mathrm{Hodge},\mathrm{nat}}(d,\lambda).
% \]
% On the other hand, by the discrete Hodge theorem,
% \[
% \ker \Delta_q^{\mathrm{Hodge},\mathrm{nat}}(d,\lambda)
% \cong
% H_q(K(d,\lambda)).
% \]
% This proves the claim. Taking dimensions gives
% \[
% \dim E_0(d,\lambda)=\beta_q(d,\lambda).
% \]
% \end{proof}


\subsection{Zero-mode spaces and the zero-mode bundle}

We retain the notation of Section~2. Let \(A\subset \mathbb{R}\) be the set of
filtration parameters, and let \(\Lambda\) be an additional parameter space.

% \begin{definition}[Parameter-dependent filtration]
% A family of finite simplicial complexes
% \[
% \{K(d,\lambda)\}_{(d,\lambda)\in A\times\Lambda}
% \]
% is called a \emph{parameter-dependent filtration} if the following conditions
% hold:
% \begin{enumerate}
%     \item for each fixed \(\lambda\in\Lambda\), the assignment
%     \[
%     d\longmapsto K(d,\lambda)
%     \]
%     is a filtration;
%     \item each \(K(d,\lambda)\) is a finite simplicial complex on a fixed finite
%     vertex set \(V\).
% \end{enumerate}
% \end{definition}

Throughout this subsection, we fix such a parameter-dependent filtration.

\begin{definition}[Ambient chain space]
Let \(K_{\max}\) denote the maximal simplicial complex on the fixed vertex set
\(V\). For each degree \(q\), define
\[
H_q:=C_q(K_{\max}).
\]
We regard \(H_q\) as a finite-dimensional Hilbert space over \(\mathbb{R}\) or
\(\mathbb{C}\), equipped with the standard inner product for which the oriented
\(q\)-simplices of \(K_{\max}\) form an orthonormal basis.
\end{definition}

For every \((d,\lambda)\in A\times\Lambda\), the chain group
\(C_q(K(d,\lambda))\) is naturally identified with a subspace of \(H_q\):
\[
C_q(K(d,\lambda))\subset H_q.
\]

\begin{definition}[Natural combinatorial Hodge Laplacian]
For each \((d,\lambda)\in A\times\Lambda\), let
\[
\partial_q(d,\lambda):
C_q(K(d,\lambda))\longrightarrow C_{q-1}(K(d,\lambda))
\]
be the boundary operator. The \emph{natural \(q\)-th combinatorial Hodge
Laplacian} on \(K(d,\lambda)\) is the self-adjoint operator
\[
\Delta_q^{\mathrm{Hodge},\mathrm{nat}}(d,\lambda)
:
C_q(K(d,\lambda))\longrightarrow C_q(K(d,\lambda))
\]
defined by
\[
\Delta_q^{\mathrm{Hodge},\mathrm{nat}}(d,\lambda)
=
\partial_q(d,\lambda)^\ast\partial_q(d,\lambda)
+
\partial_{q+1}(d,\lambda)\partial_{q+1}(d,\lambda)^\ast,
\]
where \(\ast\) denotes the adjoint with respect to the standard inner product.
\end{definition}

The operator above is defined on the parameter-dependent space
\(C_q(K(d,\lambda))\). In order to compare Laplacians at different parameter
values, we extend it to the fixed Hilbert space \(H_q\).

\begin{definition}[Extended Hodge Laplacian]
For each \((d,\lambda)\in A\times\Lambda\), consider the orthogonal
decomposition
\[
H_q
=
C_q(K(d,\lambda))\oplus C_q(K(d,\lambda))^\perp .
\]
The \emph{extended \(q\)-th Hodge Laplacian} is the self-adjoint operator
\[
\Delta_q^{\mathrm{Hodge}}(d,\lambda):H_q\longrightarrow H_q
\]
defined by
\[
\Delta_q^{\mathrm{Hodge}}(d,\lambda)
:=
\Delta_q^{\mathrm{Hodge},\mathrm{nat}}(d,\lambda)\oplus I,
\]
where \(I\) denotes the identity operator on
\(C_q(K(d,\lambda))^\perp\).
\end{definition}

\begin{remark}
The extension by the identity is chosen so that no additional zero modes are
created on the orthogonal complement. In particular,
\[
\ker \Delta_q^{\mathrm{Hodge}}(d,\lambda)
=
\ker \Delta_q^{\mathrm{Hodge},\mathrm{nat}}(d,\lambda).
\]
\end{remark}

\begin{definition}[Zero-mode space]
For each \((d,\lambda)\in A\times\Lambda\), the \emph{\(q\)-th zero-mode space}
of \(K(d,\lambda)\) is defined by
\[
E_0(d,\lambda)
:=
\ker \Delta_q^{\mathrm{Hodge}}(d,\lambda)
\subset H_q.
\]
\end{definition}

The following proposition shows that the zero-mode space recovers the usual
simplicial homology.

\begin{proposition}
For every \((d,\lambda)\in A\times\Lambda\), there is a natural isomorphism
\[
E_0(d,\lambda)\cong H_q(K(d,\lambda)).
\]
In particular,
\[
\dim E_0(d,\lambda)=\beta_q(d,\lambda).
\]
\end{proposition}

\begin{proof}
By definition of the extended Hodge Laplacian, we have
\[
\Delta_q^{\mathrm{Hodge}}(d,\lambda)
=
\Delta_q^{\mathrm{Hodge},\mathrm{nat}}(d,\lambda)\oplus I
\]
with respect to the orthogonal decomposition
\[
H_q
=
C_q(K(d,\lambda))\oplus C_q(K(d,\lambda))^\perp .
\]
Since the identity operator on \(C_q(K(d,\lambda))^\perp\) has trivial kernel,
it follows that
\[
\ker \Delta_q^{\mathrm{Hodge}}(d,\lambda)
=
\ker \Delta_q^{\mathrm{Hodge},\mathrm{nat}}(d,\lambda).
\]
On the other hand, by the discrete Hodge theorem,
\[
\ker \Delta_q^{\mathrm{Hodge},\mathrm{nat}}(d,\lambda)
\cong
H_q(K(d,\lambda)).
\]
Therefore,
\[
E_0(d,\lambda)
=
\ker \Delta_q^{\mathrm{Hodge}}(d,\lambda)
\cong
H_q(K(d,\lambda)).
\]
Taking dimensions gives
\[
\dim E_0(d,\lambda)=\beta_q(d,\lambda).
\]
This completes the proof.
\end{proof}




% \subsection{Regular regions and smooth operator models}

% We next define the regions on which the family of zero-mode spaces can be treated as a vector bundle. The purpose of this subsection is to clarify the regularity and spectral-gap assumptions needed in order to define Berry connections, curvature, and holonomy for parameter-dependent Hodge Laplacians.

% Vietoris--Rips and \v{C}ech filtrations do not, in general, vary smoothly with respect to parameters. For example, consider a finite point cloud depending on a parameter \(\lambda\),
% \[
% X(\lambda)=\{x_1(\lambda),\dots,x_N(\lambda)\}.
% \]
% The Vietoris--Rips complex at scale \(d\) is given by
% \[
% K(d,\lambda)
% =
% \{\sigma\subset X(\lambda)\mid \operatorname{diam}_\lambda(\sigma)\le d\}.
% \]
% Thus, a simplex \(\sigma\) enters the complex at the moment when one crosses the critical hypersurface
% \[
% d=\operatorname{diam}_\lambda(\sigma).
% \]
% For this reason, the ordinary combinatorial Hodge Laplacian is not, in general, a globally smooth operator family in the parameters \((d,\lambda)\). Rather, it is an operator family that changes discretely as simplices enter or leave the complex.

% In this work, we treat this discretely varying operator as a smooth Hodge-type operator family on a common Hilbert space. This does not mean that the filtration itself is replaced by a smooth topological invariant. Rather, it means that the entrance and exit of simplices are modeled by smooth activation weights in order to define Berry connections and curvature.

% We first take a maximal complex
% \[
% K_{\max}
% \]
% containing all simplices that can appear in the parameter range under consideration, and for each degree \(q\) we define the common Hilbert space
% \[
% H_q:=C_q(K_{\max}).
% \]
% This allows us to compare all chain-level operators at different parameter values \((d,\lambda)\) on the same space \(H_q\).

% For each simplex \(\sigma\in K_{\max}\), let
% \[
% r_\sigma(\lambda)
% \]
% denote its filtration threshold. In the Vietoris--Rips case, this is typically
% \[
% r_\sigma(\lambda)=\operatorname{diam}_\lambda(\sigma).
% \]
% In the ordinary filtration, the presence or absence of the simplex \(\sigma\) is described by the discontinuous indicator
% \[
% \mathbf 1_{\{d\ge r_\sigma(\lambda)\}}.
% \]
% We replace this indicator by a smooth activation weight
% \[
% w_\sigma^\varepsilon(d,\lambda)
% =
% \rho_\varepsilon(d-r_\sigma(\lambda)),
% \]
% where \(\rho_\varepsilon\) is a smooth approximation of the step function. For example, one may take
% \[
% \rho_\varepsilon(s)
% =
% \frac{1}{1+e^{-s/\varepsilon}}.
% \]
% The sigmoid function is used here only as a simple smooth approximation of the discontinuous simplex indicator. Since it is not exactly equal to \(0\) or \(1\) for finite \(\varepsilon\), the resulting weighted operator is not, in general, the exact simplicial Hodge Laplacian. Nevertheless, with a sufficiently sharp activation profile, or with a smooth cutoff function that is exactly \(0\) and \(1\) outside a small transition region, the operator agrees with, or closely approximates, the ordinary Hodge Laplacian away from critical hypersurfaces.

% Then, if \(d\ll r_\sigma(\lambda)\), we have
% \[
% w_\sigma^\varepsilon(d,\lambda)\approx 0,
% \]
% so the simplex \(\sigma\) is almost inactive. On the other hand, if \(d\gg r_\sigma(\lambda)\), then
% \[
% w_\sigma^\varepsilon(d,\lambda)\approx 1,
% \]
% so the simplex \(\sigma\) is almost fully active. Only near the critical hypersurface does \(w_\sigma^\varepsilon\) smoothly transition from \(0\) to \(1\).

% For each degree \(q\), we define the diagonal operator
% \[
% W_q^\varepsilon(d,\lambda)
% =
% \operatorname{diag}
% \bigl(
% w_\sigma^\varepsilon(d,\lambda)
% \bigr)_{\sigma\in K_{\max}^{(q)}}.
% \]
% This operator acts on the common Hilbert space \(H_q\) and records how strongly each \(q\)-simplex direction is activated.

% Let
% \[
% \partial_q:H_q\to H_{q-1}
% \]
% be the ordinary boundary operator on the maximal complex \(K_{\max}\). We define the weighted boundary operator by
% \[
% \partial_q^\varepsilon(d,\lambda)
% =
% \bigl(W_{q-1}^\varepsilon(d,\lambda)\bigr)^{1/2}
% \partial_q
% \bigl(W_q^\varepsilon(d,\lambda)\bigr)^{1/2}.
% \]
% This formula smoothly turns the boundary relation on or off according to the activation of the \(q\)-simplex on the domain side and the \((q-1)\)-simplex on the target side. If both simplices are sufficiently active, the operator is close to the ordinary boundary operator. If either simplex is inactive, its contribution is suppressed.

% From this weighted boundary operator, we define the smooth Hodge-type operator
% \[
% \Delta_q^\varepsilon(d,\lambda)
% =
% (\partial_q^\varepsilon)^\ast\partial_q^\varepsilon
% +
% \partial_{q+1}^\varepsilon(\partial_{q+1}^\varepsilon)^\ast
% +
% \mu\bigl(I-W_q^\varepsilon(d,\lambda)\bigr).
% \]
% Here \(\mu>0\) is a fixed penalty parameter. The final term
% \[
% \mu(I-W_q^\varepsilon)
% \]
% pushes inactive \(q\)-simplex directions away from the zero-mode space. Indeed, in directions where \(w_\sigma^\varepsilon\approx 1\), the penalty almost disappears, while in directions where \(w_\sigma^\varepsilon\approx 0\), a positive energy of order \(\mu\) is assigned.

% The role of \(\mu\) is only to assign positive energy to inactive simplex directions. Thus, any fixed \(\mu>0\) is sufficient from the theoretical point of view. In the numerical experiments, unless otherwise stated, we take
% \[
% \mu=1.
% \]
% This is the natural choice that assigns a unit penalty to inactive simplex directions and is sufficient to remove those directions from the zero-mode sector. More generally, one may choose \(\mu\) according to the typical scale of the nonzero spectrum, but in this work we use the fixed value \(\mu=1\).

% This construction makes \(\Delta_q^\varepsilon(d,\lambda)\) a self-adjoint and positive semidefinite operator on the common Hilbert space \(H_q\). Indeed,
% \[
% (\partial_q^\varepsilon)^\ast\partial_q^\varepsilon
% \]
% and
% \[
% \partial_{q+1}^\varepsilon(\partial_{q+1}^\varepsilon)^\ast
% \]
% are both self-adjoint and positive semidefinite, and if \(0\le W_q^\varepsilon\le I\), then
% \[
% \mu(I-W_q^\varepsilon)
% \]
% is also self-adjoint and positive semidefinite. Moreover, if \(w_\sigma^\varepsilon(d,\lambda)\) is of class \(C^r\), then \(\Delta_q^\varepsilon(d,\lambda)\) is a \(C^r\) family of self-adjoint operators.

% This smooth ambient Hodge model is not intended to define a new topological invariant. Near critical hypersurfaces, simplices may have intermediate weights between \(0\) and \(1\), and therefore
% \[
% \ker \Delta_q^\varepsilon(d,\lambda)
% \]
% need not coincide exactly with simplicial homology at every parameter value. Rather, the zero-mode space or low-energy space of \(\Delta_q^\varepsilon\) should be interpreted as the feature space of a smooth Hodge-type operator model associated with the discrete filtration. On the other hand, away from critical hypersurfaces and for small \(\varepsilon\), each \(w_\sigma^\varepsilon\) is close to either \(0\) or \(1\), and the operator approximates the ordinary Hodge Laplacian together with a positive penalty on inactive simplex directions.

% We now define regular regions. An open set \(U\subset A\times\Lambda\) is called a regular region if the chosen Hodge-type operator family
% \[
% \Delta_q(d,\lambda)
% \]
% satisfies the following conditions on \(U\).

% First,
% \[
% (d,\lambda)\mapsto \Delta_q(d,\lambda)
% \]
% is continuous, and is of class \(C^r\) whenever derivatives are needed.

% Second, the zero-mode space
% \[
% E_0(d,\lambda):=\ker \Delta_q(d,\lambda)
% \]
% has constant dimension on \(U\).

% Third, there exists a constant \(\gamma>0\) such that, for every \((d,\lambda)\in U\),
% \[
% \operatorname{Spec}\bigl(\Delta_q(d,\lambda)\bigr)
% \subset
% \{0\}\cup[\gamma,\infty).
% \]
% In other words, the zero eigenvalue is separated from the nonzero spectrum by a uniform spectral gap.

% This spectral-gap condition ensures that the zero-mode space remains stably separated from the nonzero modes. If the smallest positive eigenvalue is denoted by
% \[
% \lambda_1^+(d,\lambda)
% :=
% \min\left(
% \operatorname{Spec}(\Delta_q(d,\lambda))
% \setminus\{0\}
% \right),
% \]
% then on a regular region we require
% \[
% \lambda_1^+(d,\lambda)\ge \gamma
% \]
% uniformly.

% Under this condition, the orthogonal projection
% \[
% P_0(d,\lambda):H_q\to H_q
% \]
% onto the zero-mode space can be written as the Riesz projection
% \[
% P_0(d,\lambda)
% =
% \frac{1}{2\pi i}
% \int_\Gamma
% \bigl(z-\Delta_q(d,\lambda)\bigr)^{-1}\,dz.
% \]
% Here \(\Gamma\) is a small closed contour enclosing the origin and no nonzero spectrum. If the operator family is of class \(C^r\), then the Riesz projection also depends on \((d,\lambda)\) in a \(C^r\) manner. Hence the map
% \[
% (d,\lambda)\mapsto P_0(d,\lambda)
% \]
% is smoothly defined on the regular region.

% It follows that over a regular region \(U\),
% \[
% E_0:=\bigsqcup_{(d,\lambda)\in U}E_0(d,\lambda)
% \]
% forms a finite-dimensional vector bundle. Indeed, each fiber is
% \[
% E_0(d,\lambda)=\operatorname{Im}P_0(d,\lambda),
% \]
% and its dimension is constant on \(U\). Since \(P_0(d,\lambda)\) varies in a \(C^r\) manner, it gives a smooth family of fixed-dimensional subspaces of the common Hilbert space \(H_q\). By the standard theory of finite-dimensional Grassmannians, such a family of subspaces defines a vector bundle over \(U\).

% We call this vector bundle
% \[
% E_0\to U
% \]
% the Hodge zero-mode bundle, or the instantaneous zero-mode bundle.

% Finally, we define the singular set. Let \(\Sigma\subset A\times\Lambda\) be the set of points at which either of the following occurs:
% \begin{enumerate}
%     \item \(\dim E_0(d,\lambda)\) is not locally constant;
%     \item the spectral gap between the zero eigenvalue and the nonzero spectrum is lost.
% \end{enumerate}
% We call \(\Sigma\) the singular set, and its complement
% \[
% (A\times\Lambda)^\circ
% :=
% (A\times\Lambda)\setminus\Sigma
% \]
% the regular part.

% Thus, the theory developed in this paper is not a globally smooth theory over the entire parameter space. At critical hypersurfaces where simplices are added or removed, at points where the zero-mode rank changes, or at points where the spectral gap closes, the smooth bundle picture need not hold. Such points are excluded as part of the singular set. The curvature and holonomy considered in this work are understood as geometric quantities of the Hodge zero-mode bundle defined on regular regions away from the singular set.









\subsection{Regular regions and smooth operator models}

In this subsection, we formulate the regions on which the zero-mode spaces of
parameter-dependent Hodge Laplacians can be treated as vector bundles. In
particular, we clarify the regularity and spectral-gap assumptions needed in
order to define Berry connections, curvature, and holonomy.

Vietoris--Rips and \v{C}ech filtrations do not, in general, vary smoothly with
respect to parameters. For example, consider a finite point cloud depending on a
parameter \(\lambda\),
\[
X(\lambda)=\{x_1(\lambda),\dots,x_N(\lambda)\}.
\]
The Vietoris--Rips complex at scale \(d\) is given by
\[
K(d,\lambda)
=
\{\sigma\subset X(\lambda)\mid \operatorname{diam}_\lambda(\sigma)\le d\}.
\]
Thus, a simplex \(\sigma\) enters the complex precisely when one crosses the
critical hypersurface
\[
d=\operatorname{diam}_\lambda(\sigma).
\]
For this reason, the ordinary combinatorial Hodge Laplacian is not, in general,
a globally smooth operator family in the parameters \((d,\lambda)\). Rather, it
is an operator family that changes discretely as simplices enter or leave the
complex.

In this work, we model this discretely varying operator by a smooth Hodge-type
operator family on a common Hilbert space. This does not mean that the filtration
itself is replaced by a smooth topological invariant. Rather, the entrance and
exit of simplices are modeled by smooth activation weights so that Berry
connections and curvature can be defined.


\begin{definition}[Filtration threshold]
For each simplex \(\sigma\in K_{\max}\), let
\[
r_\sigma(\lambda)
\]
denote the filtration threshold at which \(\sigma\) appears. In the
Vietoris--Rips case, this is typically given by
\[
r_\sigma(\lambda)=\operatorname{diam}_\lambda(\sigma).
\]
\end{definition}

In the ordinary filtration, the presence or absence of the simplex \(\sigma\) is
described by the discontinuous indicator
\[
\mathbf 1_{\{d\ge r_\sigma(\lambda)\}}.
\]

\begin{definition}[Smooth activation weight]
Let \(\rho_\varepsilon\) be a smooth approximation of the step function. 
For each simplex \(\sigma\in K_{\max}\), define the smooth activation weight
\[
w_\sigma^\varepsilon(d,\lambda)
=
\rho_\varepsilon(d-r_\sigma(\lambda)).
\]
\end{definition}

For example, one may take
\[
\rho_\varepsilon(s)
=
\frac{1}{1+e^{-s/\varepsilon}}.
\]
If \(d\ll r_\sigma(\lambda)\), one has
\[
w_\sigma^\varepsilon(d,\lambda)\approx 0,
\]
so that the simplex \(\sigma\) is almost inactive. On the other hand, if
\(d\gg r_\sigma(\lambda)\), then
\[
w_\sigma^\varepsilon(d,\lambda)\approx 1,
\]
so that the simplex \(\sigma\) is almost fully active. Only near the critical
hypersurface
\[
d=r_\sigma(\lambda)
\]
does the weight smoothly transition from \(0\) to \(1\).

\begin{remark}
The sigmoid function is used only as a simple smooth approximation of the
discontinuous simplex indicator. Since it is not exactly equal to \(0\) or \(1\)
for finite \(\varepsilon\), the resulting weighted operator is not, in general,
the exact simplicial Hodge Laplacian.

Nevertheless, with a sufficiently sharp activation profile, or with a smooth
cutoff function that is exactly \(0\) and \(1\) outside a small transition
region, the operator agrees with, or closely approximates, the ordinary Hodge
Laplian away from critical hypersurfaces.
\end{remark}

\begin{definition}[Weight operator]
For each degree \(q\), define the diagonal operator
\[
W_q^\varepsilon(d,\lambda)
=
\operatorname{diag}
\bigl(
w_\sigma^\varepsilon(d,\lambda)
\bigr)_{\sigma\in K_{\max}^{(q)}}.
\]
This is a self-adjoint operator on the common Hilbert space \(H_q\), and records
the activation strength of each \(q\)-simplex direction.
\end{definition}

\begin{definition}[Weighted boundary operator]
Let
\[
\partial_q:H_q\to H_{q-1}
\]
be the ordinary boundary operator on the maximal complex \(K_{\max}\). The
weighted boundary operator is defined by
\[
\partial_q^\varepsilon(d,\lambda)
=
\bigl(W_{q-1}^\varepsilon(d,\lambda)\bigr)^{1/2}
\partial_q
\bigl(W_q^\varepsilon(d,\lambda)\bigr)^{1/2}.
\]
\end{definition}

This definition smoothly turns the boundary relation on or off according to the
activation of the \(q\)-simplex on the domain side and the \((q-1)\)-simplex on
the target side. If both simplices are sufficiently active, the operator is close
to the ordinary boundary operator. If either simplex is inactive, its
contribution is suppressed.

\begin{definition}[Smooth Hodge-type operator]
Fix a parameter \(\mu>0\). The smooth Hodge-type operator in degree \(q\) is
defined by
\[
\Delta_q^\varepsilon(d,\lambda)
=
(\partial_q^\varepsilon)^\ast\partial_q^\varepsilon
+
\partial_{q+1}^\varepsilon(\partial_{q+1}^\varepsilon)^\ast
+
\mu\bigl(I-W_q^\varepsilon(d,\lambda)\bigr).
\]
Here \(\mu\) is a penalty parameter assigning positive energy to inactive
\(q\)-simplex directions.
\end{definition}

The final term
\[
\mu(I-W_q^\varepsilon)
\]
pushes inactive \(q\)-simplex directions away from the zero-mode space. Indeed,
in directions where \(w_\sigma^\varepsilon\approx 1\), the penalty almost
disappears, while in directions where \(w_\sigma^\varepsilon\approx 0\), a
positive energy of order \(\mu\) is assigned.

\begin{remark}[Choice of the penalty parameter]
The role of \(\mu\) is only to assign positive energy to inactive simplex
directions. Thus, from the theoretical point of view, any fixed \(\mu>0\) is
sufficient.

In the numerical experiments, unless otherwise stated, we take
\[
\mu=1.
\]
This is the natural choice that assigns a unit penalty to inactive simplex
directions and is sufficient to remove those directions from the zero-mode
sector. More generally, one may choose \(\mu\) according to the typical scale of
the nonzero spectrum, but in this work we use the fixed value \(\mu=1\).
\end{remark}

\begin{proposition}[Basic properties of the smooth Hodge-type operator]
The operator
\[
\Delta_q^\varepsilon(d,\lambda)
\]
defined above is self-adjoint and positive semidefinite on the common Hilbert
space \(H_q\).

Moreover, if each weight \(w_\sigma^\varepsilon(d,\lambda)\) is of class \(C^r\)
in \((d,\lambda)\), then
\[
(d,\lambda)\longmapsto \Delta_q^\varepsilon(d,\lambda)
\]
is a \(C^r\) family of self-adjoint operators.
\end{proposition}

\begin{proof}
First,
\[
(\partial_q^\varepsilon)^\ast\partial_q^\varepsilon
\]
is of the form \(A^\ast A\), with \(A=\partial_q^\varepsilon\). Hence it is
self-adjoint and positive semidefinite. Similarly,
\[
\partial_{q+1}^\varepsilon(\partial_{q+1}^\varepsilon)^\ast
\]
is of the form \(BB^\ast\), and is therefore self-adjoint and positive
semidefinite.

Furthermore, if
\[
0\le w_\sigma^\varepsilon(d,\lambda)\le 1
\]
for all \(\sigma\), then the diagonal operator \(W_q^\varepsilon(d,\lambda)\)
satisfies
\[
0\le W_q^\varepsilon(d,\lambda)\le I.
\]
Hence
\[
I-W_q^\varepsilon(d,\lambda)
\]
is self-adjoint and positive semidefinite. Since \(\mu>0\), the operator
\[
\mu(I-W_q^\varepsilon(d,\lambda))
\]
is also self-adjoint and positive semidefinite.

Therefore \(\Delta_q^\varepsilon(d,\lambda)\), being the sum of these three
terms, is self-adjoint and positive semidefinite.

If \(w_\sigma^\varepsilon(d,\lambda)\) depends on \((d,\lambda)\) in a \(C^r\)
manner, then so does \(W_q^\varepsilon(d,\lambda)\). Moreover, since
\(W_q^\varepsilon(d,\lambda)\) is diagonal with nonnegative entries, its square
root also depends \(C^r\)-smoothly on \((d,\lambda)\), under the chosen smooth
activation model. Consequently,
\[
\partial_q^\varepsilon(d,\lambda)
\]
is a \(C^r\) family of operators, and hence so is
\[
\Delta_q^\varepsilon(d,\lambda).
\]
\end{proof}

\begin{remark}
The smooth ambient Hodge model is not intended to define a new topological
invariant. Near critical hypersurfaces, simplices may have intermediate weights
between \(0\) and \(1\), and therefore
\[
\ker \Delta_q^\varepsilon(d,\lambda)
\]
need not coincide exactly with simplicial homology at every parameter value.

Rather, the zero-mode space, or low-energy space, of
\(\Delta_q^\varepsilon(d,\lambda)\) should be interpreted as the feature space of
a smooth Hodge-type operator model associated with the discrete filtration. On
the other hand, away from critical hypersurfaces and for sufficiently small
\(\varepsilon\), each \(w_\sigma^\varepsilon\) is close to either \(0\) or \(1\),
and the operator approximates the ordinary Hodge Laplacian together with a
positive penalty on inactive simplex directions.
\end{remark}

\begin{definition}[Regular region]
Let \(U\subset A\times\Lambda\) be an open set, and let
\[
\Delta_q(d,\lambda)
\]
be the chosen Hodge-type operator family in degree \(q\). We say that \(U\) is a
\emph{regular region} for \(\Delta_q\) if the following conditions hold:

\begin{enumerate}
    \item The map
    \[
    (d,\lambda)\longmapsto \Delta_q(d,\lambda)
    \]
    is continuous on \(U\), and is of class \(C^r\) whenever derivatives are
    needed.

    \item The zero-mode space
    \[
    E_0(d,\lambda):=\ker \Delta_q(d,\lambda)
    \]
    has constant dimension on \(U\).

    \item There exists a constant \(\gamma>0\) such that, for every
    \((d,\lambda)\in U\),
    \[
    \operatorname{Spec}\bigl(\Delta_q(d,\lambda)\bigr)
    \subset
    \{0\}\cup[\gamma,\infty).
    \]
\end{enumerate}
\end{definition}

Thus, on a regular region, the zero eigenvalue is separated from the nonzero
spectrum by a uniform spectral gap.

\begin{remark}[Spectral gap condition]
The spectral-gap condition ensures that the zero-mode space remains stably
separated from the nonzero modes.

If the smallest positive eigenvalue is denoted by
\[
\lambda_1^+(d,\lambda)
:=
\min\left(
\operatorname{Spec}(\Delta_q(d,\lambda))
\setminus\{0\}
\right),
\]
then on a regular region we require
\[
\lambda_1^+(d,\lambda)\ge \gamma
\]
uniformly.
\end{remark}

\begin{proposition}[Riesz projection onto the zero-mode space]
Let \(U\subset A\times\Lambda\) be a connected regular region. Then the orthogonal
projection
\[
P_0(d,\lambda):H_q\to H_q
\]
onto the zero-mode space is given by the Riesz projection
\[
P_0(d,\lambda)
=
\frac{1}{2\pi i}
\int_\Gamma
\bigl(z-\Delta_q(d,\lambda)\bigr)^{-1}\,dz.
\]
Here \(\Gamma\) is a small closed contour enclosing the origin and no nonzero
spectrum.

Moreover, if \(\Delta_q(d,\lambda)\) is a \(C^r\) operator family, then
\[
(d,\lambda)\longmapsto P_0(d,\lambda)
\]
is also of class \(C^r\).
\end{proposition}

\begin{proof}
By the definition of a regular region, the zero eigenvalue is separated from the
nonzero spectrum by a uniform positive distance. Therefore, one can choose a
closed contour \(\Gamma\), uniformly over \(U\), which encloses the origin and no
nonzero spectrum.

By the standard spectral decomposition for finite-dimensional operators, or
equivalently by the Riesz functional calculus,
\[
\frac{1}{2\pi i}
\int_\Gamma
\bigl(z-\Delta_q(d,\lambda)\bigr)^{-1}\,dz
\]
is the spectral projection onto the generalized eigenspace corresponding to the
eigenvalues enclosed by \(\Gamma\).

Since \(\Delta_q(d,\lambda)\) is self-adjoint and the only eigenvalue enclosed by
\(\Gamma\) is \(0\), this projection is precisely the orthogonal projection onto
\[
E_0(d,\lambda)=\ker \Delta_q(d,\lambda).
\]
This proves the formula.

If \(\Delta_q(d,\lambda)\) is of class \(C^r\), then for each \(z\in\Gamma\) the
resolvent
\[
\bigl(z-\Delta_q(d,\lambda)\bigr)^{-1}
\]
also depends \(C^r\)-smoothly on \((d,\lambda)\). Since the contour \(\Gamma\) is
fixed, integration over \(\Gamma\) preserves \(C^r\)-regularity. Hence
\[
P_0(d,\lambda)
\]
is a \(C^r\) family of projections.
\end{proof}

\begin{proposition}[Zero-mode bundle]
Let \(U\subset A\times\Lambda\) be a connected regular region. Then
\[
E_0
:=
\bigsqcup_{(d,\lambda)\in U}E_0(d,\lambda)
\]
forms a finite-dimensional vector bundle over \(U\). Its fiber at
\((d,\lambda)\in U\) is
\[
E_0(d,\lambda)=\operatorname{Im}P_0(d,\lambda).
\]
\end{proposition}

\begin{proof}
By the definition of a regular region,
\[
\dim E_0(d,\lambda)
\]
is constant on \(U\). Moreover, by the preceding proposition, the projection
\[
P_0(d,\lambda)
\]
depends continuously, or \(C^r\)-smoothly, on \((d,\lambda)\).

Therefore,
\[
\operatorname{Im}P_0(d,\lambda)
\]
defines a continuous, or \(C^r\), family of fixed-dimensional subspaces of the
fixed finite-dimensional Hilbert space \(H_q\).
% Equivalently, this gives a continuous, or \(C^r\), map
% \[
% U\to \operatorname{Gr}(k,H_q),
% \qquad
% (d,\lambda)\mapsto E_0(d,\lambda),
% \]
% where
% \[
% k=\dim E_0(d,\lambda).
% \]
% By the standard theory of finite-dimensional Grassmannians, such a family of
% subspaces defines a finite-dimensional vector bundle over \(U\). 
Hence
\[
E_0
=
\bigsqcup_{(d,\lambda)\in U}E_0(d,\lambda)
\]
is a vector bundle over \(U\).
\end{proof}

\begin{definition}[Hodge zero-mode bundle]
Let \(U\subset A\times\Lambda\) be a connected regular region. The vector bundle
\[
E_0\to U
\]
constructed above is called the \emph{\(q\)-th Hodge zero-mode bundle}, or the
\emph{instantaneous zero-mode bundle}, over \(U\).
\end{definition}

\begin{definition}[Singular set and regular part]
Let \(\Sigma\subset A\times\Lambda\) be the set of points at which at least one
of the following occurs:
\begin{enumerate}
    \item \(\dim E_0(d,\lambda)\) is not locally constant;
    \item the spectral gap between the zero eigenvalue and the nonzero spectrum
    is lost.
\end{enumerate}
We call \(\Sigma\) the \emph{singular set}. Its complement
\[
(A\times\Lambda)^\circ
:=
(A\times\Lambda)\setminus\Sigma
\]
is called the \emph{regular part}.
\end{definition}


\begin{remark}
The geometric theory developed in this paper is not a globally smooth theory
over the entire parameter space. At critical hypersurfaces where the spectral gap closes, the smooth vector-bundle picture need not hold.
Such points are excluded as part of the singular set \(\Sigma\). The Berry
connections, curvature, and holonomy considered in this work are understood as
geometric quantities of the Hodge zero-mode bundle defined on regular regions
away from the singular set.

Strictly speaking, if the regular region \(U\) is not connected, the
zero-mode multiplicity may be constant only on each connected component and
may differ from one component to another. In that case, \(E_0\) should not be
regarded as a single vector bundle of fixed rank over all of \(U\), but rather as
a collection of vector bundles over the connected components of \(U\). In the
following discussion, we always work on one connected component on which
\(\dim E_0(d,\lambda)=m\) is fixed. For simplicity, we continue to call this object
the zero-mode bundle.
\end{remark}
















\subsection{Natural connection and Berry connection}

Let \(U\subset A\times\Lambda\) be a regular region, and consider the zero-mode bundle
\[
E_0\to U.
\]
For each \((d,\lambda)\in U\), the fiber \(E_0(d,\lambda)\) is a subspace of the common Hilbert space \(H_q\), and \(P_0(d,\lambda)\) is the orthogonal projection onto it. Hence \(P_0\) induces a natural connection on \(E_0\).

\begin{definition}
For a local section \(s\in\Gamma(U,E_0)\), define
\[
\nabla s:=P_0(ds).
\]
\end{definition}

\begin{proposition}
The operator \(\nabla\) defines a connection on the vector bundle \(E_0\to U\).
\end{proposition}

\begin{proof}
Linearity of \(\nabla\) follows from the definition. Let \(f\in C^\infty(U)\) and \(s\in\Gamma(U,E_0)\). Then
\[
\nabla(fs)
=
P_0(d(fs))
=
P_0(df\otimes s+f\,ds).
\]
Since \(s\) is a section of \(E_0\), we have \(P_0s=s\) pointwise. Therefore
\[
P_0(df\otimes s)=df\otimes s,
\qquad
P_0(f\,ds)=fP_0(ds)=f\nabla s.
\]
Thus
\[
\nabla(fs)=df\otimes s+f\nabla s.
\]
This is the Leibniz rule, so \(\nabla\) is a connection.
\end{proof}

Let \(\operatorname{rank}E_0=m\), and choose a local orthonormal frame
\[
\psi_1(d,\lambda),\dots,\psi_m(d,\lambda)
\]
on \(U\). Arrange these vectors as columns into the matrix
\[
\Psi(d,\lambda)
:=
(\psi_1(d,\lambda),\dots,\psi_m(d,\lambda)).
\]
Then
\[
\Psi^\ast\Psi=I_m,
\qquad
P_0=\Psi\Psi^\ast.
\]

Every local section \(s\) can be written uniquely as
\[
s=\Psi\xi,
\]
where \(\xi:U\to\mathbb C^m\) is a coefficient vector. In the real case, \(\mathbb C^m\) is replaced by \(\mathbb R^m\). Then
\[
\nabla s
=
\Psi\bigl(d\xi+\Psi^\ast d\Psi\,\xi\bigr).
\]
Thus the connection \(1\)-form in the frame \(\Psi\) is
\[
A:=\Psi^\ast d\Psi.
\]

\begin{definition}
The matrix-valued \(1\)-form
\[
A=\Psi^\ast d\Psi
\]
is called the Berry connection on the zero-mode bundle.
\end{definition}

Differentiating \(\Psi^\ast\Psi=I_m\), we obtain
\[
d\Psi^\ast\,\Psi+\Psi^\ast d\Psi=0.
\]
Hence
\[
A^\ast=-A.
\]
In particular, in the complex case, \(A\) is a \(\mathfrak u(m)\)-valued \(1\)-form. In the real case, it is an \(\mathfrak o(m)\)-valued \(1\)-form.

If the local orthonormal frame is changed by
\[
\Psi'=\Psi g,
\]
then in the complex case \(g:U\to U(m)\), while in the real case \(g:U\to O(m)\). The corresponding connection \(1\)-form \(A'\) satisfies
\[
A'=g^{-1}Ag+g^{-1}dg.
\]
Thus the local expression of the Berry connection is gauge-dependent.

\subsection{Curvature and holonomy}

We now define the curvature and holonomy associated with the Berry connection.

\begin{definition}
The curvature of the Berry connection \(A\) is defined by
\[
F:=dA+A\wedge A.
\]
\end{definition}

If \(U\) is a two-dimensional parameter space with local coordinates \((d,t)\), then
\[
A=A_d\,dd+A_t\,dt,
\]
and
\[
F=F_{dt}\,dd\wedge dt,
\]
where
\[
F_{dt}
=
\partial_dA_t-\partial_tA_d+[A_d,A_t].
\]

The curvature can also be expressed solely in terms of the projection \(P_0\).

\begin{proposition}
The curvature satisfies
\[
F=P_0[dP_0,dP_0]P_0.
\]
\end{proposition}

\begin{proof}
This follows by a standard computation using \(P_0=\Psi\Psi^\ast\) and \(A=\Psi^\ast d\Psi\).
\end{proof}

This projection formula is particularly important in the present work. The connection form \(A\) depends on the choice of local frame, whereas the projection \(P_0\) is intrinsically defined as an operator on the common Hilbert space. Thus, by writing the curvature as
\[
F=P_0[dP_0,dP_0]P_0,
\]
we obtain a gauge-independent description of the noncommutativity of local zero-mode transport.

Next, let \(\gamma:[0,1]\to U\) be a piecewise \(C^1\) curve. The parallel-transport equation in a local frame is
\[
\frac{d}{ds}\xi(s)
+
A(\dot\gamma(s))\xi(s)
=
0.
\]
Its solution defines a linear map
\[
\Pi_\gamma:E_0(\gamma(0))\to E_0(\gamma(1)).
\]

\begin{definition}
If \(\gamma\) is a closed curve, that is, if \(\gamma(0)=\gamma(1)=x\), then
\[
U_\gamma:=\Pi_\gamma:E_0(x)\to E_0(x)
\]
is called the holonomy along \(\gamma\).
\end{definition}

In a local frame, the holonomy is written as
\[
U_\gamma
=
\mathcal{P}\exp\left(-\int_\gamma A\right),
\]
where \(\mathcal{P}\) denotes path ordering.

Under a gauge transformation \(\Psi'=\Psi g\), the matrix representation of the holonomy transforms by conjugation:
\[
U_\gamma'
=
g(\gamma(0))^{-1}U_\gamma g(\gamma(0)).
\]
Therefore,
\[
\operatorname{tr}(U_\gamma),
\qquad
\det(U_\gamma),
\qquad
\operatorname{Spec}(U_\gamma)
\]
are gauge-invariant.

\subsection{Stability on regular regions}

Finally, we state the local stability of the zero-mode projection and the curvature on regular regions. To obtain an intrinsic stability statement that does not depend on gauge choices, we do not directly compare local connection forms \(A\) or holonomy matrices. Instead, we focus on \(P_0\), which can be compared as an operator on the common Hilbert space, and on the curvature expressed by the projection formula
\[
F=P_0[dP_0,dP_0]P_0.
\]

Let \(U\subset A\times\Lambda\) be a regular region, and let
\[
L(d,\lambda),
\qquad
\widetilde L(d,\lambda)
\]
be two self-adjoint operator families on \(H_q\). Here we regard them as common-Hilbert-space realizations of ordinary Hodge Laplacians, or as perturbations of such realizations. Let
\[
P_0,\qquad \widetilde P_0
\]
denote the corresponding zero-mode projections. Define the corresponding curvatures by
\[
F:=P_0[dP_0,dP_0]P_0,
\qquad
\widetilde F:=
\widetilde P_0[d\widetilde P_0,d\widetilde P_0]\widetilde P_0.
\]


\begin{theorem}[Quantitative stability on regular regions]
Assume that:
\begin{enumerate}
\item \(L\) and \(\widetilde L\) are of class \(C^2\) on \(U\) and have the same constant zero-mode multiplicity \(m\);
\item there exists a common spectral-gap constant \(\gamma>0\) such that, for every \(x\in U\),
\[
\operatorname{Spec}(L(x)),
\operatorname{Spec}(\widetilde L(x))
\subset
\{0\}\cup[\gamma,\infty);
\]
\item \(L\) and \(\widetilde L\) have uniformly bounded \(C^2\) norms on \(U\).
\end{enumerate}
Then there exist constants \(C_0,C_1,C_2>0\) such that
\[
\|P_0-\widetilde P_0\|_{C^0(U)}
\le
C_0\|L-\widetilde L\|_{C^0(U)},
\]
\[
\|dP_0-d\widetilde P_0\|_{C^0(U)}
\le
C_1\|L-\widetilde L\|_{C^1(U)},
\]
and
\[
\|F-\widetilde F\|_{C^0(U)}
\le
C_2\|L-\widetilde L\|_{C^2(U)}.
\]
In particular, on a regular region with a uniform spectral gap, the zero-mode projection is Lipschitz stable with respect to \(C^0\)-perturbations of the operator family, and the curvature is Lipschitz stable with respect to \(C^2\)-perturbations of the operator family.
\end{theorem}

\begin{proof}[Proof sketch]
The zero-mode projection is expressed as the Riesz projection
\[
P_0
=
\frac{1}{2\pi i}
\int_\Gamma
(z-L)^{-1}\,dz.
\]
Hence the estimate for \(P_0-\widetilde P_0\) follows from the resolvent identity. Next, differentiating the Riesz projection formula gives the estimate for \(dP_0-d\widetilde P_0\). Finally, using the projection formula for curvature,
\[
F=P_0[dP_0,dP_0]P_0,
\]
the estimate for \(F-\widetilde F\) follows from the estimates for \(P_0\) and \(dP_0\).
\end{proof}


\end{document}
